% ==========================================
% SECCIÓN 3: INVESTIGACIÓN BIBLIOGRÁFICA
% ==========================================

\section{INVESTIGACIÓN BIBLIOGRÁFICA Y ESTADO DEL ARTE}

\subsection{Trabajos Previos y Antecedentes de Automatización Académica}

En el ámbito informático de las Instituciones de Educación Superior (IES) del Ecuador y de la región, la sustitución progresiva de herramientas ofimáticas aisladas por plataformas web centralizadas representa una tendencia consolidada para fortalecer la gobernanza, seguridad y trazabilidad de los datos institucionales \cite{kendall_systems}, \cite{pressman_eng}. A nivel nacional, existen antecedentes concretos de modernización y automatización de procesos académico-administrativos que fundamentan la pertinencia técnica y la necesidad operativa de la presente investigación:

\begin{itemize}[leftmargin=*]
    \item \textbf{Gestión y Automatización Documental de Becas Universitarias:} Un precedente directo de sustancial relevancia lo constituye el trabajo de Romero y Medina (UNESUM, 2026) \cite{romero_unesum_becas}, titulado \textit{«Diseño de un sistema informático para la automatización de la gestión documental en el área de becas de la Universidad Estatal del Sur de Manabí»}. Dicho proyecto abordó de manera frontal la problemática del manejo físico y disperso en hojas de cálculo, demostrando que la implementación de una arquitectura web modular con control de acceso basado en roles (RBAC) y almacenamiento estructurado en la nube no solo erradica la duplicidad operacional y el extravío de información, sino que salvaguarda la integridad de los expedientes de los beneficiarios en estricto apego a las disposiciones de la Ley Orgánica de Educación Superior (LOES) \cite{loes} y la Ley Orgánica de Protección de Datos Personales (LOPDP) \cite{lopdp_ecuador}.
    
    \item \textbf{Gestión Automatizada y Trazabilidad de Trámites Académicos:} En la Universidad Técnica de Ambato (UTA, 2021), Barrionuevo \cite{barrionuevo_uta_tramites} desarrolló la investigación \textit{«Sistema automatizado para la gestión de trámites académicos en la secretaría de la FISEI»}. Los hallazgos empíricos de este estudio evidenciaron que reemplazar los canales de recepción física y las comunicaciones manuales por una plataforma web centralizada para la radicación de solicitudes en línea con seguimiento de estados en tiempo real reduce drásticamente los cuellos de botella burocráticos, minimiza los tiempos de respuesta y eleva la satisfacción usuaria en la comunidad académica.
    
    \item \textbf{Sistemas Estatales de Seguimiento Contractual y Devengamiento de Becarios:} En el plano gubernamental ecuatoriano, la Secretaría de Educación Superior, Ciencia, Tecnología e Innovación (SENESCYT) opera la plataforma oficial \textbf{PUSAK} \cite{senescyt_pusak}. Este sistema funciona como el referente institucional a nivel nacional para la administración de expedientes, registro de avances académicos, control de adendas o prórrogas y liquidación contractual de becarios nacionales e internacionales. El modelo funcional de PUSAK corrobora empíricamente la necesidad de contar con algoritmos parametrizados capaces de calcular y supervisar de forma automatizada los plazos de devengamiento y compromisos legales de los becarios.
    
    \item \textbf{Control Centralizado de Expedientes Docentes y Normalización Relacional:} En la Universidad Metropolitana del Ecuador (UMET) \cite{umet_expedientes}, se han implementado plataformas web para la centralización, auditoría y control de expedientes docentes. Estas experiencias institucionales confirman que la sustitución de archivos planos de Microsoft Excel por esquemas relacionales normalizados en tercera forma normal (3FN) sobre motores robustos como PostgreSQL o MySQL erradica hasta en un 95\% las inconsistencias, redundancias y errores humanos de registro presentes en la administración ofimática dispersa \cite{elmasri}, \cite{silberschatz}.
\end{itemize}

A nivel internacional, la literatura especializada en ingeniería de software corrobora que la adopción de arquitecturas web modulares desacopladas orientadas a servicios API RESTful centralizados transforma la interoperabilidad entre dependencias universitarias centrales y facultades académicas \cite{fielding_rest}, \cite{newman_microservices}. Este enfoque desacoplado viabiliza el intercambio ágil y seguro de datos de filiación y estado contractual, garantizando que los módulos administrativos y los portales de autogestión de usuarios operen con alta cohesión, bajo acoplamiento y óptima mantenibilidad \cite{fowler}, \cite{sommerville_eng}.

\subsection{Identificación de la Brecha Tecnológica Institucional}

A pesar de los avances reportados en la literatura y en desarrollos previos, la revisión bibliográfica e institucional revela una brecha tecnológica crítica en la administración central de la UCE. Actualmente no existe una plataforma web diseñada expresamente para resolver la lógica contractual única que rige a la Unidad de Gestión de Doctorados. Las soluciones genéricas existentes no contemplan el algoritmo dinámico requerido para calcular prórrogas, adendas y periodos de devengamiento docente de conformidad con la Ley Orgánica de Educación Superior (LOES) y el Reglamento de Régimen Académico de la UCE \cite{loes}, \cite{reglamento_uce}.

Adicionalmente, los sistemas de gestión académica universitarios suelen carecer de interfaces estandarizadas y desacopladas para el intercambio seguro de expedientes docentes con la administración central. Esta desarticulación dificulta la sincronización ágil de datos de filiación con los servicios institucionales de la Dirección de Tecnologías de la Información y Comunicación (DTICs). Por consiguiente, el presente proyecto integrador resuelve esta brecha mediante una plataforma web desacoplada de alto rendimiento (Next.js 16, Astro, Express.js, Prisma ORM y PostgreSQL) que combina el cálculo automatizado de devengamientos, alertas automáticas de vencimiento, generación de certificados digitales, atención de solicitudes mediante tickets y contratos de interfaz desacoplados (\textit{mock services}) listos para una migración directa a los servicios de la DTICs.
